\section{Introdução}

Endodontia e prótese dentária são, historicamente, especialidades que operam em sequência: o tratamento endodôntico restaura a sanidade do complexo pulpar e periapical, criando as condições biológicas necessárias para que a reabilitação protética devolva função e estética ao elemento dental. Essa interdependência clínica, no entanto, sempre foi gerenciada de forma fragmentada, com cada especialidade apoiando-se em exames e registros independentes, frequentemente incomunicáveis entre si. A transição para uma odontologia de precisão exige a superação dessa fragmentação por meio de infraestruturas computacionais unificadas \cite{techsoc2025precision}.

\destaque{A consolidação de \textit{Digital Twins} (Gêmeos Digitais) na prática odontológica introduz uma ruptura paradigmática nesse fluxo de trabalho, permitindo uma integração contínua e bidirecional de dados.}

Pela primeira vez, é possível construir uma réplica virtual unificada do dente -- abrangendo desde a morfologia interna dos canais radiculares até a geometria externa da coroa e as propriedades mecânicas dos materiais restauradores -- que acompanha o paciente ao longo de todo o ciclo de tratamento. Esse modelo único permite que endodontistas e protesistas planejem suas intervenções de forma coordenada. Simulações do comportamento biomecânico do conjunto dente--pino--coroa--osso podem ser conduzidas antes mesmo da primeira instrumentação, antecipando falhas estruturais e otimizando a escolha dos materiais \cite{khoshfekr2025digital}. Adicionalmente, a convergência digital atua como um catalisador para a criação de protocolos que reduzem a variabilidade técnica entre os operadores \cite{mdpi2025convergence}.

Este capítulo examina em profundidade as aplicações emergentes de gêmeos digitais na endodontia e na prótese dentária. Na primeira seção, discutem-se os avanços no diagnóstico endodôntico tridimensional, no monitoramento preditivo de instrumentos de NiTi, na endodontia guiada e na simulação de procedimentos regenerativos. Na segunda seção, a atenção volta-se para as próteses fixas, sobre implante, removíveis e as reabilitações orais complexas, enfatizando as inovações em simulação mastigatória e oclusal \cite{roy2025applications}. A terceira seção aborda a integração endo-protética per se, seguida pelos desafios metodológicos e perspectivas futuras.

Embora a endodontia e a prótese apresentem graus distintos de maturidade tecnológica, ambas compartilham um horizonte epistemológico comum: a transição para uma odontologia preditiva, personalizada e balizada por modelos computacionais matematicamente validados. O desenvolvimento de arcabouços de integração, como o padrão \textit{Functional Mock-up Interface} e os sistemas multiagentes \cite{infante2024fmi, infante2025distributed}, fornece a base para que essas especialidades deixem de ser compartimentos estanques e passem a operar como componentes de um ecossistema holístico.
